\documentclass[11pt,a4paper]{article}

% =========================================================
% Packages
% =========================================================
\usepackage[utf8]{inputenc}
\usepackage[T1]{fontenc}
\usepackage{lmodern}
\usepackage[a4paper,margin=2.2cm]{geometry}
\usepackage{xcolor}
\usepackage{tcolorbox}
\tcbuselibrary{breakable}
\usepackage{enumitem}
\usepackage{hyperref}

% =========================================================
% Style
% =========================================================
\definecolor{studyblue}{HTML}{2563EB}
\definecolor{studygray}{HTML}{F3F4F6}

\setlength{\parindent}{0pt}
\setlist[itemize]{leftmargin=1.4em, itemsep=2pt, topsep=4pt}

\newtcolorbox{definition}[1]{
    breakable,
    title=\textbf{#1},
    colback=studyblue!4,
    colframe=studyblue!65!black,
    boxrule=0.8pt,
    arc=2mm,
    left=3mm,
    right=3mm,
    top=2mm,
    bottom=2mm,
    before skip=8pt,
    after skip=8pt
}

\newtcolorbox{studylist}[1]{
    breakable,
    title=\textbf{#1},
    colback=studygray,
    colframe=black!40,
    boxrule=0.6pt,
    arc=2mm,
    left=3mm,
    right=3mm,
    top=2mm,
    bottom=2mm,
    before skip=8pt,
    after skip=8pt
}

% =========================================================
% Document
% =========================================================
\title{Software Design -- Klassediagram}
\author{Dries Van den Brande, Stef De Bie, Lars Verheyden}
\date{}

\begin{document}

\maketitle

\section{Klassediagram}

% =========================================================
\subsection{Basisconcepten}
% =========================================================

\begin{definition}{Klasse}
Een verzameling van objecten met overeenkomstige eigenschappen.

\textbf{Kennen:}
\begin{itemize}
    \item Een klassebeschrijving bevat de naam van de klasse en haar eigenschappen.
    \item Klassen worden door de modelleur bepaald, bijvoorbeeld via classificatie.
\end{itemize}
\end{definition}

\begin{definition}{Object / instantie}
Een object heeft een zelfstandig bestaan, in werkelijkheid of conceptueel.

\textbf{Kennen:}
\begin{itemize}
    \item \textbf{Identiteit}: elk object is uniek.
    \item \textbf{Toestand}: de data van het object.
    \item \textbf{Gedrag}: de operaties van het object.
    \item \textbf{Encapsulation}: interne details worden verborgen.
    \item \textbf{Abstraction}: het object wordt via zijn interface benaderd.
    \item \textbf{Instantie} is een andere naam voor object, vooral wanneer men wil aangeven tot welke klasse het object behoort.
\end{itemize}
\end{definition}

\begin{definition}{Attribuut}
Informatie die door een object beheerd wordt.

\textbf{Kennen:}
\begin{itemize}
    \item De waarden van alle attributen samen vormen de toestand van een object.
    \item Twee objecten kunnen dezelfde toestand hebben en toch verschillende objecten zijn, omdat elk object een eigen identiteit heeft.
    \item Een attribuut heeft een bepaald type.
    \item Voorbeelden bij een persoon: naam, geboortedatum en gewicht.
    \item In UML kan een attribuut een multipliciteit hebben.
\end{itemize}
\end{definition}

\begin{definition}{Operatie}
Een operatie beschrijft een service die door een object geleverd wordt.

\textbf{Kennen:}
\begin{itemize}
    \item Een operatie moet behoren tot de verantwoordelijkheden van het object.
    \item De verzameling operaties vormt het gedrag van het object.
    \item Een operatie kan argumenten hebben.
    \item Een operatie kan een resultaat opleveren.
    \item Objecten communiceren door boodschappen te sturen tussen zender en ontvanger.
\end{itemize}
\end{definition}

\begin{definition}{Associatie}
Een structurele relatie tussen twee klassen.

\textbf{Kennen:}
\begin{itemize}
    \item Een instantie van de ene klasse is gedurende een groot deel van haar bestaan verbonden met een instantie van de andere klasse.
    \item Tussen twee klassen kunnen meerdere associaties bestaan.
    \item De naam van de associatie staat langs de associatielijn.
    \item De leesrichting kan aangeduid worden.
    \item Een \textbf{link} is een instantie van een associatie.
\end{itemize}
\end{definition}

\begin{definition}{Rol}
De rol beschrijft een associatie-einde.

\textbf{Kennen:}
\begin{itemize}
    \item Een object kan in verschillende associaties verschillende rollen spelen.
    \item Wanneer de rolnaam wordt weggelaten, geldt de klassenaam als rolnaam.
\end{itemize}
\end{definition}

\begin{definition}{Multipliciteit}
Geeft aan hoeveel instanties aan een associatie-einde of attribuut kunnen voorkomen.

\textbf{Kennen:}
\begin{itemize}
    \item Multipliciteit kan met symbolen of expliciete cijfers worden weergegeven.
    \item Multipliciteit kan aan beide uiteinden van een associatie staan.
    \item Bij een attribuut wordt de multipliciteit tussen vierkante haken geplaatst.
\end{itemize}
\end{definition}

\begin{definition}{Generalisatie}
Het classificeren van klassen volgens een \textbf{IS-A}-relatie.

\textbf{Kennen:}
\begin{itemize}
    \item Zoek klassen met gemeenschappelijke kenmerken.
    \item \textbf{Superklasse} = generalisatie = basisklasse.
    \item \textbf{Subklasse} = specialisatie = afgeleide klasse.
\end{itemize}
\end{definition}

\begin{definition}{Overerving}
Iedere subklasse erft de eigenschappen van haar superklasse.

\textbf{Kennen:}
\begin{itemize}
    \item Attributen en operaties van de superklasse worden geërfd.
    \item Overerving volgt de IS-A-relatie.
\end{itemize}
\end{definition}

\begin{definition}{Liskov's substitutieprincipe}
Op iedere plaats waar een instantie van een klasse gebruikt kan worden, kan een instantie van een subklasse gebruikt worden.

\textbf{Kennen:}
\begin{itemize}
    \item Het principe hoort bij de IS-A-relatie.
\end{itemize}
\end{definition}

\begin{definition}{Commentaar}
Extra uitleg in een UML-diagram.

\textbf{Kennen:}
\begin{itemize}
    \item Commentaar kan in een notebox geplaatst worden.
\end{itemize}
\end{definition}

\begin{definition}{Klassediagram}
Een diagram dat de klassen in een systeem toont.

\textbf{Kennen:}
\begin{itemize}
    \item Bevat klassen.
    \item Kan attributen bevatten.
    \item Kan operaties bevatten.
    \item Kan associaties bevatten.
\end{itemize}
\end{definition}

\begin{definition}{Objectdiagram}
Een afbeelding van objecten op een bepaald tijdstip.

\textbf{Kennen:}
\begin{itemize}
    \item De objecten zijn gecreëerd volgens de structuur van het klassediagram.
    \item Een objectdiagram toont instanties van het klassediagram.
\end{itemize}
\end{definition}

% =========================================================
\subsection{Geavanceerde concepten}
% =========================================================

\begin{definition}{Afgeleid attribuut}
Een attribuut waarvan de waarde afgeleid wordt uit andere informatie.

\textbf{Kennen:}
\begin{itemize}
    \item De waarde kan afgeleid worden uit andere attributen, associaties of objecten.
    \item De waarde hoeft niet noodzakelijk opgeslagen te worden.
    \item Het is informatie waar een object om gevraagd kan worden.
    \item Voorbeeld: \texttt{/leeftijd}.
\end{itemize}
\end{definition}

\begin{definition}{Zichtbaarheid}
Geeft aan in hoeverre een attribuut of operatie door de omgeving van een object gekend is.

\textbf{Kennen:}
\begin{itemize}
    \item \texttt{-} = private
    \item \texttt{+} = public
    \item \texttt{\#} = protected
    \item \texttt{\textasciitilde} = package
    \item Inkapseling probeert zoveel mogelijk interne informatie te verbergen.
\end{itemize}
\end{definition}

\begin{definition}{Bag}
Een meervoudige verzameling waarbij volgorde niet belangrijk is en dubbele waarden toegelaten zijn.
\end{definition}

\begin{definition}{Sequence}
Een meervoudige verzameling met een vaste volgorde waarin dubbele waarden toegelaten zijn.
\end{definition}

\begin{definition}{Set}
Een meervoudige verzameling zonder vaste volgorde waarin iedere waarde uniek is.
\end{definition}

\begin{definition}{OrderedSet}
Een meervoudige verzameling met een vaste volgorde waarin iedere waarde uniek is.
\end{definition}

\begin{definition}{Klasseattribuut}
Een attribuut dat een eigenschap van de klasse beschrijft en niet van iedere individuele instantie.

\textbf{Kennen:}
\begin{itemize}
    \item Heeft slechts één waarde, onafhankelijk van het aantal instanties van de klasse.
\end{itemize}
\end{definition}

\begin{definition}{Klasseoperatie}
Een service die door de klasse zelf geleverd wordt en niet door een individuele instantie.

\textbf{Kennen:}
\begin{itemize}
    \item Een constructor is een voorbeeld dat in de slides bij klasseoperaties wordt vermeld.
\end{itemize}
\end{definition}

\begin{definition}{Afgeleide associatie}
Een associatie die uit andere informatie kan worden afgeleid.

\textbf{Kennen:}
\begin{itemize}
    \item Is redundant.
    \item Moet spaarzaam gebruikt worden.
\end{itemize}
\end{definition}

\begin{definition}{Richting van een associatie}
Geeft aan welke klasse de andere klasse kent.

\textbf{Kennen:}
\begin{itemize}
    \item Een gewone associatie impliceert in principe dat de objecten aan beide uiteinden elkaar kennen.
    \item Een richting kan aangeduid worden om deze navigatie te beperken.
    \item Dit kan consistentieproblemen helpen vermijden.
\end{itemize}
\end{definition}

\begin{definition}{Associatieklasse}
Een klasse die gekoppeld is aan een associatie.

\textbf{Kennen:}
\begin{itemize}
    \item De link tussen twee instanties is een instantie van de associatieklasse.
    \item Wordt gebruikt wanneer de associatie zelf attributen heeft.
    \item Wordt gebruikt wanneer de associatie zelf operaties heeft.
    \item Wordt gebruikt wanneer de associatie associaties met andere klassen heeft.
\end{itemize}
\end{definition}

\begin{definition}{Compositie}
Een speciaal soort associatie waarbij één of meer klassen onderdeel zijn van een andere klasse.

\textbf{Kennen:}
\begin{itemize}
    \item Een onderdeel mag in een compositierelatie slechts tot één geheel behoren.
    \item De levensduur van het deelobject is kleiner dan of gelijk aan de levensduur van het gehele object.
    \item Er is een \textbf{lifetime dependency}.
\end{itemize}
\end{definition}

\begin{definition}{Abstracte klasse}
Een klasse waarvan geen instanties kunnen bestaan.

\textbf{Kennen:}
\begin{itemize}
    \item Is superklasse van andere klassen.
    \item Dient om een gezamenlijke interface te definiëren.
    \item Concrete subklassen kunnen wel instanties hebben.
\end{itemize}
\end{definition}

\begin{definition}{Abstracte operatie}
Een operatie die in een abstracte klasse gedeclareerd wordt maar daar niet volledig gedefinieerd hoeft te zijn.
\end{definition}

\begin{definition}{Meervoudige overerving}
Een situatie waarbij een klasse meer dan één superklasse heeft.
\end{definition}

\begin{definition}{Herhaalde overerving}
Een vorm van overerving die kan ontstaan binnen een hiërarchie met meerdere overervingspaden.

\textbf{Kennen:}
\begin{itemize}
    \item De slides vermelden dat meervoudige en herhaalde overerving niet in moderne programmeertalen gebruikt worden.
\end{itemize}
\end{definition}

\begin{definition}{Enumeratie}
Een type waarvoor een vast aantal mogelijke waarden gedefinieerd is.
\end{definition}

\begin{definition}{Package}
Een groepering van modelelementen.

\textbf{Kennen:}
\begin{itemize}
    \item Is een generiek mechanisme om een softwaresysteem op te delen.
    \item Ieder modelelement behoort tot één package.
    \item Associaties kunnen voorkomen tussen klassen uit verschillende packages.
    \item Een package is een verzameling klassen die samenwerken om aan een bepaalde verantwoordelijkheid te voldoen.
    \item Packages kunnen genest zijn.
    \item Een klasse uit een andere package kan met de notatie \texttt{PACKAGE::Klasse} aangeduid worden.
    \item Sinds UML 2 bestaat ook \textit{package merge}.
\end{itemize}
\end{definition}

\begin{definition}{Geparametriseerde klasse / template}
Een klasse waarvan bepaalde zaken afhankelijk zijn van één of meer parameters.

\textbf{Kennen:}
\begin{itemize}
    \item Door waarden aan de parameters te geven ontstaat een klasse waarvan instanties gemaakt kunnen worden.
    \item Een template beschrijft een familie van klassen.
    \item Wanneer een parameter een type aanduidt, kan aangegeven worden aan welk ander type dat type moet voldoen.
\end{itemize}
\end{definition}

\begin{definition}{Constraint}
Een beperking op één of meer elementen in een klassediagram.

\textbf{Kennen:}
\begin{itemize}
    \item Constraints kunnen met OCL uitgedrukt worden.
    \item Voorbeelden uit de slides zijn leeftijdsvoorwaarden en XOR-constraints.
\end{itemize}
\end{definition}

\begin{definition}{OCL}
Object Constraint Language.

\textbf{Kennen:}
\begin{itemize}
    \item Wordt gebruikt om bedrijfsregels en beperkingen bij een model uit te drukken.
    \item Een invariant kan in de vorm \texttt{context ... inv: ...} geschreven worden.
\end{itemize}
\end{definition}

% =========================================================
\subsection{Werkwijze voor een klassediagram}
% =========================================================

\begin{studylist}{Stappenplan}
\begin{enumerate}
    \item Identificeer alle mogelijke kandidaatklassen.
    \item Selecteer de klassen uit de lijst van kandidaten.
    \item Maak eventueel een modeldictionary.
    \item Identificeer associaties.
    \item Identificeer attributen.
    \item Identificeer operaties.
    \item Generaliseer met behulp van overerving.
    \item Voeg eventueel OCL-constraints toe.
    \item Groepeer eventueel in packages.
    \item Itereer over de gedane stappen.
\end{enumerate}
\end{studylist}

\begin{definition}{Kandidaatklasse identificeren}
Zoek zonder beperkingen zoveel mogelijk mogelijke klassen.

\textbf{Kennen:}
\begin{itemize}
    \item Methode 1: onderstreep zelfstandige naamwoorden in de tekst.
    \item Methode 2: organiseer een brainstormsessie.
    \item Betrek toekomstige gebruikers, domeinexperts en ontwikkelaars.
    \item Tijdens deze stap: geen kritiek.
    \item Niet iedere term wordt uiteindelijk een klasse; termen kunnen wel in de modeldictionary terechtkomen.
\end{itemize}
\end{definition}

\begin{definition}{Klassen selecteren}
Bepaal welke kandidaatklassen effectief in het model opgenomen worden.

\textbf{Vragen:}
\begin{itemize}
    \item Heeft de kandidaat een zelfstandige rol in het probleemdomein?
    \item Willen we iets van deze kandidaat weten?
    \item Moet de kandidaat iets doen?
    \item Heeft de kandidaat een duidelijke verantwoordelijkheid in het systeem?
\end{itemize}

\textbf{Mogelijke redenen om geen klasse te maken:}
\begin{itemize}
    \item Redundant.
    \item Irrelevant voor het probleemdomein.
    \item Te vaag.
    \item Niet kwantificeerbaar.
    \item Is eigenlijk een attribuut.
    \item Is eigenlijk een operatie.
    \item Is eigenlijk een rol.
\end{itemize}
\end{definition}

\begin{definition}{Modeldictionary}
Een verzameling waarin de termen van het model worden opgenomen.

\textbf{Kennen:}
\begin{itemize}
    \item Kan klassen, attributen, operaties en andere termen bevatten.
\end{itemize}
\end{definition}

\begin{definition}{Associaties identificeren}
Zoek relaties tussen paren van klassen.

\textbf{Kennen:}
\begin{itemize}
    \item Zoek kandidaatassociaties.
    \item Elimineer tijdelijke acties of relaties die eerder dependencies zijn.
    \item Controleer of afgeleide associaties echt nodig zijn.
    \item Koppel terug naar gebruikers en domeinexperts.
    \item Voeg rollen en multipliciteiten toe.
    \item Verwijder irrelevante associaties.
\end{itemize}
\end{definition}

\begin{definition}{Attributen identificeren}
Breid de klassen uit met relevante informatie die door hun objecten beheerd wordt.

\textbf{Vragen:}
\begin{itemize}
    \item Wat zijn de dingen waar een object van weet?
    \item Welke informatie kunnen we aan het object vragen?
\end{itemize}

\textbf{Opletten voor:}
\begin{itemize}
    \item Een attribuut kan geen objectreferentie vervangen; daarvoor gebruik je een associatie.
    \item Identifiers zijn vaak kunstmatige attributen.
    \item Neem alleen interne waarden op die voor de klant relevant zijn.
    \item Vermijd te veel detail.
    \item Denk ook aan klasseattributen.
    \item Voeg multipliciteit toe.
    \item Typen worden volgens de slides later in het ontwerp toegevoegd.
\end{itemize}
\end{definition}

\begin{definition}{Operaties identificeren}
Bepaal welk gedrag bij de verantwoordelijkheid van een klasse hoort.

\textbf{Vragen:}
\begin{itemize}
    \item Wat wil ik dat het object voor mij doet?
    \item Wat is de verantwoordelijkheid van het object?
\end{itemize}
\end{definition}

\begin{definition}{Generaliseren met overerving}
Bekijk of bestaande klassen gemeenschappelijke eigenschappen hebben die in een nieuwe superklasse geplaatst kunnen worden.

\textbf{Kennen:}
\begin{itemize}
    \item Zoek identieke attributen, operaties en associaties.
    \item Controleer of ze ook werkelijk dezelfde betekenis hebben.
    \item Generaliseren is een vorm van abstraheren.
    \item De slides waarschuwen dat te veel generalisatie de flexibiliteit kan verminderen.
\end{itemize}
\end{definition}

\begin{definition}{Bedrijfsregels toevoegen}
Druk business rules en beperkingen uit met constraints, bijvoorbeeld in OCL.

\textbf{Kennen:}
\begin{itemize}
    \item Formuleer vragen binnen het probleemdomein.
    \item Wanneer belangrijke vragen niet met het model beantwoord kunnen worden, kan het model te beperkt zijn.
\end{itemize}
\end{definition}

\begin{definition}{Klassen groeperen in packages}
Deel een groot klassediagram logisch of praktisch op in packages.

\textbf{Kennen:}
\begin{itemize}
    \item Binnen UML is de gekozen verdeling vrij.
\end{itemize}
\end{definition}

\begin{definition}{Itereren}
Het maken van een klassediagram is geen strikt sequentieel proces.

\textbf{Kennen:}
\begin{itemize}
    \item Controleer op inconsistenties en onvolledigheid.
    \item Een verantwoordelijkheid of operatie zonder geschikte klasse kan wijzen op een ontbrekend modelelement.
    \item Een constraint die niet beschreven kan worden met het bestaande model kan wijzen op een ontbrekende associatie, klasse, attribuut of operatie.
    \item Pas het diagram aan en doorloop eerdere stappen opnieuw.
\end{itemize}
\end{definition}

\end{document}
